\section{Simultaneous equivariant trivialization}

\begin{theorem}[Scalar homogeneous cocycle is trivial]\label{thm:gauge}
The function \(c\) in \eqref{eq:cocycle} is a chronological one-cocycle.
The same gauge
\begin{equation}\label{eq:gauge}
  T(x,v)=(x,\phi(x)^{-1}v)
\end{equation}
simultaneously trivializes the twisted \(\Q^\times\)-descent and the full
idele-scaling lift.  Hence its transformation-groupoid \(H^1\)-class and
its equivariant line-bundle class are zero.
\end{theorem}

\begin{proof}
The cocycle identity is the telescoping equality
\[
  c(ab,x)=\frac{\phi(abx)}{\phi(x)}
  =c(a,bx)c(b,x).
\]
Twist either the rational or idele action by
\(g\cdot(x,v)=(gx,c(g,x)v)\).  Then
\[
  T(gx,c(g,x)v)
  =(gx,\phi(gx)^{-1}c(g,x)v)
  =(gx,\phi(x)^{-1}v),
\]
which is the untwisted action in the new coordinate.  Because the formula
is identical for rational descent and idele scaling, their compatibility
square is trivialized as well.
\end{proof}

This statement is stronger than checking a few local vacua.  It says that
every gauge-invariant representation, descent, or determinant functor sees
the trivial scalar class.

\begin{corollary}[Prime repetitions]\label{cor:prime}
For every prime \(p\) and \(r\ge1\), the H\'enon scalar holonomy around the
scaling loop of length \(r\log p\) equals one.
\end{corollary}

\begin{proof}
Along the successive scales, the exponent is
\[
  \sum_{j=0}^{r-1}2(p^3-1)p^{3j}x^3
  =2(p^{3r}-1)x^3.
\]
The rational endpoint identification contributes the inverse gauge
\(\Pz(x)-\Pz(p^rx)=-2(p^{3r}-1)x^3\).  Their product is one.  Equivalently,
this is the closed-loop form of the telescoping proof of
\cref{thm:gauge}.
\end{proof}

The clock and the inherited scaling Euler factor may remain meaningful,
but they are not generated by a nontrivial scalar H\'enon holonomy.  Any
opposite claim must exhibit exactly where functoriality or gauge
admissibility is broken.
